% !TeX root = ../main.tex

% ==================================================
% 2.11. 메서드의 종류 — 인스턴스 메서드, 클래스 메서드, 정적 메서드
% ==================================================

\subsection{메서드의 종류 - 인스턴스 메서드, 클래스 메서드, 정적 메서드}

2.5에서 인스턴스 메서드가 클래스에 저장된 함수 객체이며, 인스턴스를 통해 접근하면 바운드 메서드(Bound Method)가 된다는 사실을 학습했다.

2.7에서 디스크립터 프로토콜을, 2.8과 2.9에서 상속과 MRO를 살펴보았다.

이번 2.11에서는 이 개념들을 연결하여 Python의 세 가지 메서드를 비교해 볼 것이다.

\begin{dhksmry}
    인스턴스 메서드 - self

    클래스 메서드 - cls, @classmethod

    정적 메서드 - @staticmethod
\end{dhksmry}

이번 학습의 핵심 질문은 다음과 같다.

\begin{dhksmry}[centered]
세 메서드는 모두 클래스 안에서 정의되는데,

왜 어떤 메서드는 인스턴스를 자동으로 전달받고,

어떤 메서드는 클래스를 전달받으며,

어떤 메서드는 아무것도 자동으로 전달받지 않을까?
\end{dhksmry}

이 차이를 이해하면 단순히 문법을 구분하는 수준을 넘어, Python이 함수 객체를 어떻게 바인딩하는지까지 설명할 수 있다.

\subsubsection{세 가지 메서드의 기본 구조}

먼저 다음 클래스를 살펴보자.

\begin{dhkcode}
class Person:
    species = "Human"

    def __init__(self, name):
        self.name = name

    def introduce(self):
        return f"My name is {self.name}"

    @classmethod
    def get_species(cls):
        return cls.species

    @staticmethod
    def is_valid_name(name):
        return isinstance(name, str) and len(name) > 0
\end{dhkcode}

이 클래스에는 세 가지 메서드가 정의되어 있다.

\begin{dhkcode}
p = Person("Kim")

print(p.introduce())
print(Person.get_species())
print(Person.is_valid_name("Lee"))
\end{dhkcode}

실행 결과:

\begin{dhkzsh}
My name is Kim
Human
True
\end{dhkzsh}

겉으로 보면 모두 메서드 호출이다.

하지만 내부적으로 전달되는 첫 번째 인수가 서로 다르다.

\begin{dhktable}
|구분|	정의 방식|	자동 전달 대상|	주된 역할|
|:---|:---|:---|:---|
|인스턴스 메서드|	def method(self)|	인스턴스|	인스턴스 상태 사용|
|클래스 메서드|	@classmethod|	클래스 객체|	클래스 상태 및 클래스 기반 생성|
|정적 메서드|	@staticmethod|	없음|	클래스와 관련된 독립적 기능|
\end{dhktable}

여기서 중요한 점은 \zsh{self}와 \zsh{cls}가 Python의 예약어가 아니라는 사실이다.

둘 다 일반적인 매개변수 이름이다.

실제로 중요한 것은 메서드 접근 시 Python이 어떤 객체를 자동으로 바인딩하는가이다.

\subsubsection{인스턴스 메서드 - \zsh{self}는 어떻게 전달되는가?}

다음 코드를 보자.

\begin{dhkcode}
class Person:
    def __init__(self, name):
        self.name = name

    def introduce(self):
        return self.name

p = Person("Kim")

print(p.introduce())
\end{dhkcode}

실행 결과:

\begin{dhkzsh}
Kim
\end{dhkzsh}

2.5에서 배운 것처럼 다음 두 표현식은 일반적인 메서드 호출에서 같은 효과를 가진다.

\begin{dhkcode}
p.introduce()

Person.introduce(p)
\end{dhkcode}

왜 가능할까?

클래스에 정의된 일반 함수 객체는 디스크립터 프로토콜의 \zsh{__get__()}을 제공하기 때문이다.

\begin{dhknote}
\begin{dhksmry}[centered]
클래스 네임스페이스 확인
\end{dhksmry}
\end{dhknote}

\begin{dhkcode}
print(Person.__dict__["introduce"])
\end{dhkcode}

출력 예시:

\begin{dhkzsh}
<function Person.introduce at 0x...>
\end{dhkzsh}

\zsh{Person.__dict__}에 저장된 것은 함수 객체다.

반면 인스턴스를 통해 접근하면 다음과 같다.

\begin{dhkcode}
method = p.introduce

print(method.__self__ is p)
print(method.__func__ is Person.introduce)
\end{dhkcode}

실행 결과:

\begin{dhkzsh}
True
True
\end{dhkzsh}

여기서

\begin{dhksmry}
    __self__는 바인딩된 인스턴스

    __func__는 원본 함수 객체
\end{dhksmry}

를 나타낸다.

\begin{dhkgraph}[uniform]
node n1 (1,1) :: introduce -> 함수
node n2 (1,3) :: name -> "Kim"
node n3 (2,2) [aceent] :: __self__->p // __func__-> 함수
group g1(n1) :: Person 클래스
group g2(n2) :: p 인스턴스
group g3(n3) :: Bound Method 
merge g1,g2 -> g3

caption :: p.introduce에 접근하면 인스턴스와 함수 객체가 연결된 바운드 메서드를 얻는다.
\end{dhkgraph}

따라서 인스턴스 메서드의 핵심은 단순히 \zsh{self}라는 이름을 사용한다는 것이 아니다.

함수 객체가 인스턴스에 바인딩되어 호출될 때 그 인스턴스가 첫 번째 인수로 전달된다는 것이 핵심이다.

\subsubsection{클래스 메서드 - \zsh{@classmethod}는 무엇을 바꾸는가?}

이번에는 클래스 메서드를 살펴보자.

\begin{dhkcode}
class Person:
    species = "Human"

    @classmethod
    def get_species(cls):
        return cls.species
\end{dhkcode}

다음처럼 호출할 수 있다.

\begin{dhkcode}
print(Person.get_species())
\end{dhkcode}

실행 결과:

\begin{dhkzsh}
Human
\end{dhkzsh}

여기서 \zsh{cls}는 \zsh{Person} 클래스 객체를 참조한다.

따라서 다음과 같은 효과가 있다.

\begin{dhkcode}
Person.get_species()

Person.get_species.__func__(Person)
\end{dhkcode}

두 번째 표현식은 바운드 메서드에 저장된 원본 함수를 명시적으로 호출한 형태다.

\begin{dhknote}
\begin{dhksmry}[centered]
cls는 어디에서 오는가?
\end{dhksmry}
\end{dhknote}

\zsh{@classmethod}는 일반 함수 객체를 \zsh{classmethod} 디스크립터 객체로 감싸는 데코레이터다.

2.5에서 학습했던 데코레이터 변환식을 적용하면 다음과 같이 이해할 수 있다.

\begin{dhkcode}
class Person:
    def get_species(cls):
        return cls.species

    get_species = classmethod(get_species)
\end{dhkcode}

이는 다음 정의와 대응한다.

\begin{dhkcode}
class Person:
    @classmethod
    def get_species(cls):
        return cls.species
\end{dhkcode}

중요한 점은 \zsh{classmethod}가 원본 함수 객체의 코드를 변경하는 것이 아니라, 클래스를 바인딩하는 방식으로 접근할 수 있게 하는 디스크립터 객체를 생성한다는 사실이다.

\begin{dhknote}
\begin{dhksmry}[centered]
클래스 네임스페이스 확인
\end{dhksmry}
\end{dhknote}

\begin{dhkcode}
class Person:
    @classmethod
    def get_species(cls):
        return cls.__name__

print(type(Person.__dict__["get_species"]))
\end{dhkcode}

실행 결과:

\begin{dhkzsh}
<class 'classmethod'>
\end{dhkzsh}

일반 인스턴스 메서드와는 다르다.

일반 메서드는 클래스 네임스페이스에 함수 객체가 직접 저장된다.

반면 \zsh{@classmethod}를 사용하면 \zsh{classmethod} 객체가 저장된다.

\begin{dhknote}
\begin{dhksmry}[centered]
바운드 메서드 확인
\end{dhksmry}
\end{dhknote}

\begin{dhkcode}
method = Person.get_species

print(method.__self__ is Person)
print(method.__func__ is Person.__dict__["get_species"].__func__)
\end{dhkcode}

실행 결과:

\begin{dhkzsh}
True
True
\end{dhkzsh}

여기서 가장 중요한 부분은 다음이다.

\begin{dhkcode}
method.__self__ is Person
\end{dhkcode}

인스턴스 메서드에서는 \zsh{__self__}가 인스턴스였지만, 클래스 메서드에서는 클래스 객체다.

즉, 두 메서드는 모두 바운드 메서드가 될 수 있지만 바인딩 대상이 다르다는 것이다.

\subsubsection{클래스 메서드를 인스턴스에서 호출하면?}

다음 코드를 보자.

\begin{dhkcode}
class Person:
    species = "Human"

    @classmethod
    def get_species(cls):
        print(cls)
        return cls.species

p = Person()

p.get_species()
\end{dhkcode}

출력:

\begin{dhkzsh}
<class '__main__.Person'>
\end{dhkzsh}

여기서 \zsh{p.get_species()}를 호출했는데도 \zsh{cls}는 \zsh{p}가 아니다.

\zsh{Person} 클래스 객체다.

확인해 보자.

\begin{dhkcode}
print(p.get_species.__self__ is Person)
\end{dhkcode}

실행 결과:

\begin{dhkzsh}
True
\end{dhkzsh}

즉, 클래스 메서드는 클래스 객체로 접근하든 인스턴스로 접근하든 클래스를 바인딩한다.

이것은 인스턴스 메서드와 중요한 차이다.

\begin{dhknote}
\begin{dhksmry}[centered]
비교
\end{dhksmry}
\end{dhknote}

\begin{dhkcode}
class Person:
    def instance_method(self):
        return self

    @classmethod
    def class_method(cls):
        return cls

p = Person()

print(p.instance_method() is p)
print(p.class_method() is Person)
\end{dhkcode}

실행 결과:

\begin{dhkzsh}
True
True
\end{dhkzsh}

\subsubsection{클래스 메서드와 상속 - \zsh{cls}는 어떤 클래스를 참조하는가?}

이제 2.8의 상속과 연결해 보자.

\begin{dhkcode}
class Animal:
    species = "Animal"

    @classmethod
    def identify(cls):
        return cls.species


class Dog(Animal):
    species = "Dog"
\end{dhkcode}

다음 코드를 실행하면 어떻게 될까?

\begin{dhkcode}
print(Animal.identify())
print(Dog.identify())
\end{dhkcode}

실행 결과:

\begin{dhkzsh}
Animal
Dog
\end{dhkzsh}

여기서 중요한 질문이 있다.

\zsh{identify()}는 \zsh{Animal} 클래스에 정의되어 있다.

그런데 \zsh{Dog.identify()}를 호출하면 왜 \zsh{Dog}가 출력될까?

클래스 메서드의 \zsh{cls}는 메서드가 정의된 클래스로 고정되는 것이 아니라, 접근에 사용된 클래스에 따라 바인딩되기 때문이다.

\begin{dhkcode}
print(Animal.identify.__self__ is Animal)
print(Dog.identify.__self__ is Dog)
\end{dhkcode}

실행 결과:

\begin{dhkzsh}
True
True
\end{dhkzsh}

이 구조는 2.8에서 학습한 인스턴스 메서드의 \zsh{self}와 비슷하다.

부모 클래스에 정의된 메서드를 자식 클래스가 사용할 수 있고, 그때 실제 바인딩 대상이 달라지는 것이다.

\begin{dhknote}
\begin{dhksmry}[centered]
중요한 비교
\end{dhksmry}
\end{dhknote}

\begin{dhktable}
|구분|	인스턴스 메서드|	클래스 메서드|
|:---|:---|:---|
|바인딩 대상|	실제 인스턴스|	접근 대상 클래스|
|관례적 첫 인수|	self|	cls|
|상속 시|	자식 인스턴스를 self로 받을 수 있음|	자식 클래스를 cls로 받을 수 있음|
\end{dhktable}

이 차이가 클래스 메서드를 활용하는 중요한 이유다.

\subsubsection{클래스 메서드의 대표적 용도 - 대체 생성자}

클래스 메서드가 특히 유용한 경우는 

\zsh{대체 생성자(Alternative Constructor)}를 정의할 때다.

먼저 다음 코드를 보자.

\begin{dhkcode}
class Person:
    def __init__(self, name, age):
        self.name = name
        self.age = age
\end{dhkcode}

일반적인 인스턴스 생성은 다음과 같다.

\begin{dhkcode}
p = Person("Kim", 30)
\end{dhkcode}

그런데 다음과 같은 문자열에서 객체를 생성하고 싶다고 하자.

\begin{dhkcode}
data = "Kim,30"
\end{dhkcode}

이를 위해 클래스 메서드를 사용할 수 있다.

\begin{dhkcode}
class Person:
    def __init__(self, name, age):
        self.name = name
        self.age = age

    @classmethod
    def from_string(cls, data):
        name, age = data.split(",")

        return cls(name, int(age))
\end{dhkcode}

이제 다음처럼 사용할 수 있다.

\begin{dhkcode}
p = Person.from_string("Kim,30")

print(p.name)
print(p.age)
\end{dhkcode}

실행 결과:

\begin{dhkzsh}
Kim
30
\end{dhkzsh}

여기서 다음 표현식에 주목하라.

\begin{dhkcode}
return cls(name, int(age))
\end{dhkcode}

\zsh{cls}는 클래스 객체다.

따라서 \zsh{cls(...)}는 해당 클래스 객체를 호출하여 인스턴스를 생성하는 동작이다.

2.6에서 학습했던 클래스 호출 메커니즘과 연결된다.

\begin{dhknote}
\begin{dhksmry}[centered]
왜 Person(...) 대신 cls(...)를 사용할까?
\end{dhksmry}
\end{dhknote}

상속 관계를 고려하면 차이가 명확해진다.

\begin{dhkcode}
class Student(Person):
    pass
\end{dhkcode}

다음 코드를 실행해 보자.

\begin{dhkcode}
student = Student.from_string("Lee,20")

print(type(student))
\end{dhkcode}

실행 결과:

\begin{dhkzsh}
<class '__main__.Student'>
\end{dhkzsh}

왜 \zsh{Student} 인스턴스가 생성되었을까?

\zsh{Student.from_string()}으로 접근했기 때문에 

클래스 메서드의 \zsh{cls}가 \zsh{Student}로 바인딩되었기 때문이다.

따라서

\begin{dhkcode}
return cls(name, int(age))
\end{dhkcode}

는 이 호출에서 사실상 다음과 같은 효과를 가진다.

\begin{dhkcode}
return Student(name, int(age))
\end{dhkcode}

반면 클래스 메서드 내부에서 다음처럼 작성했다면

\begin{dhkcode}
return Person(name, int(age))
\end{dhkcode}

항상 \zsh{Person} 인스턴스가 생성된다.

자식 클래스의 타입을 유지할 수 없을 것이다.

이것이 클래스 메서드가 대체 생성자에 적합한 중요한 이유다.

단, 자식 클래스가 생성자 인수를 다르게 정의했다면 \zsh{cls(...)} 호출이 실패할 수 있으므로 상속 계층의 생성자 규약도 호환되어야 한다.

\subsubsection{정적 메서드 - \zsh{@staticmethod}는 무엇을 하는가?}

이번에는 정적 메서드를 살펴보자.

\begin{dhkcode}
class Person:
    @staticmethod
    def is_valid_age(age):
        return isinstance(age, int) and age >= 0
\end{dhkcode}

다음처럼 호출할 수 있다.

\begin{dhkcode}
print(Person.is_valid_age(30))
print(Person.is_valid_age(-5))
\end{dhkcode}

실행 결과:

\begin{dhkzsh}
True
False
\end{dhkzsh}

여기서 \zsh{is_valid_age()}에는 \zsh{self}도 없고 \zsh{cls}도 없다.

왜 그럴까?

\zsh{@staticmethod}는 인스턴스나 클래스를 자동으로 바인딩하지 않기 대문이다.

즉, 정적 메서드는 클래스의 네임스페이스에 속해 있지만, 호출할 때 자동으로 추가되는 첫 번째 인수가 없다.

\begin{dhknote}
\begin{dhksmry}[centered]
데코레이터 변환식
\end{dhksmry}
\end{dhknote}

다음 코드는

\begin{dhkcode}
class Person:
    @staticmethod
    def is_valid_age(age):
        return age >= 0
\end{dhkcode}

개념적으로 다음과 대응한다.

\begin{dhkcode}
class Person:
    def is_valid_age(age):
        return age >= 0

    is_valid_age = staticmethod(is_valid_age)
\end{dhkcode}

클래스 네임스페이스에는 \zsh{staticmethod} 객체가 저장된다.

\begin{dhkcode}
print(type(Person.__dict__["is_valid_age"]))
\end{dhkcode}

실행 결과:

\begin{dhkzsh}
<class 'staticmethod'>
\end{dhkzsh}

하지만 클래스나 인스턴스를 통해 접근하면 일반적으로 감싸진 함수 객체를 얻는다.

\begin{dhkcode}
p = Person()

print(Person.is_valid_age is p.is_valid_age)
\end{dhkcode}

실행 결과:

\begin{dhkzsh}
True
\end{dhkzsh}

둘 다 같은 함수 객체를 반환할 수 있다.

인스턴스나 클래스에 바인딩하는 과정이 없기 때문이다.

\begin{dhknote}
\begin{dhksmry}[centered]
staticmethod는 왜 필요한가?
\end{dhksmry}
\end{dhknote}

다음 함수를 생각해 보자.

\begin{dhkcode}
def is_valid_age(age):
    return isinstance(age, int) and age >= 0
\end{dhkcode}

이 함수는 독립적으로 존재해도 문제가 없다.

하지만 \zsh{Person}과 관련된 검증 기능이므로 클래스의 네임스페이스 안에 배치하고 싶을 수 있다.

\begin{dhkcode}
class Person:
    @staticmethod
    def is_valid_age(age):
        return isinstance(age, int) and age >= 0
\end{dhkcode}

이렇게 하면 다음과 같이 기능을 묶을 수 있다.

\begin{dhkcode}
Person.is_valid_age(30)
\end{dhkcode}

정적 메서드는 클래스와 개념적으로 관련된 함수를 클래스의 네임스페이스 안에 배치하면서, 인스턴스나 클래스를 자동으로 전달하지 않도록 하는 수단이라고 이해하면 된다.

다만 함수가 클래스와 특별한 관련이 없다면 굳이 정적 메서드로 만들 필요는 없다.

모듈 수준의 일반 함수로 정의하는 것이 더 명확할 수도 있다.

\subsubsection{세 메서드의 바인딩 구조를 직접 비교하기}

이번에는 지금까지의 내용을 하나의 코드로 비교해 보자.

\begin{dhkcode}
class Example:
    def instance_method(self):
        return self

    @classmethod
    def class_method(cls):
        return cls

    @staticmethod
    def static_method():
        return "Hello"

obj = Example()
\end{dhkcode}

각 메서드에 접근했을 때 어떤 객체를 얻는지 확인해 보자.

\begin{dhkcode}
print(type(obj.instance_method))
print(type(obj.class_method))
print(type(obj.static_method))
\end{dhkcode}

실행 결과:

\begin{dhkzsh}
<class 'method'>
<class 'method'>
<class 'function'>
\end{dhkzsh}

여기서 중요한 결과가 나타났다.

인스턴스 메서드와 클래스 메서드는 모두 바운드 메서드 객체로 접근된다.

반면 정적 메서드는 일반 함수 객체로 접근된다.

\begin{dhknote}
\begin{dhksmry}[centered]
__self__ 비교
\end{dhksmry}
\end{dhknote}

\begin{dhkcode}
print(obj.instance_method.__self__ is obj)
print(obj.class_method.__self__ is Example)
\end{dhkcode}

실행 결과:

\begin{dhkzsh}
True
True
\end{dhkzsh}

정적 메서드는 일반 함수 객체이므로 바운드 메서드의 \zsh{__self__}를 갖지 않는다.

\begin{dhkcode}
print(hasattr(obj.static_method, "__self__"))
\end{dhkcode}

실행 결과:

\begin{dhkzsh}
False
\end{dhkzsh}

세 가지를 그림으로 정리하면 다음과 같다.

\begin{dhkgraph}[uniform]
node n1 (1,1) :: 인스턴스 메서드 // obj.instance_method // \da // Bound Method -> __self__ = obj
node n2 (2,1) :: 클래스 메서드 // obj.class_method // \da // Bound Method -> __self__ = Example
node n3 (3,1) :: 정적 메서드 // obj.static_method // \da // Function -> 자동 바인딩 없음

caption :: 세 메서드의 핵심 차이는 함수 본문이 아니라 속성 접근 시 적용되는 바인딩 방식이다.
\end{dhkgraph}

\subsubsection{디스크립터 프로토콜로 세 메서드를 설명하기}

이제 2.7에서 학습한 디스크립터 개념을 연결해 보자.

세 메서드의 차이는 클래스 네임스페이스에 저장되는 객체와 그 객체의 \zsh{__get__()} 동작으로 설명할 수 있다.

다음 클래스를 정의하자.

\begin{dhkcode}
class Example:
    def instance_method(self):
        pass

    @classmethod
    def class_method(cls):
        pass

    @staticmethod
    def static_method():
        pass
\end{dhkcode}

클래스 딕셔너리의 객체 타입을 확인해 보자.

\begin{dhkcode}
for name in (
    "instance_method",
    "class_method",
    "static_method"
):
    obj = Example.__dict__[name]
    print(name, type(obj))
\end{dhkcode}

실행 결과:

\begin{dhkzsh}
instance_method <class 'function'>
class_method <class 'classmethod'>
static_method <class 'staticmethod'>
\end{dhkzsh}

이제 구조를 정리할 수 있다.

\begin{dhktable}
|메서드|	클래스 __dict__에 저장되는 객체|	__get__() 결과|
|:---|:---|:---|
|인스턴스 메서드|	function|	인스턴스에 바인딩된 메서드|
|클래스 메서드|	classmethod|	클래스에 바인딩된 메서드|
|정적 메서드|	staticmethod|	감싸진 함수 객체|
\end{dhktable}

실제로 \zsh{__get__()}을 직접 호출해 볼 수도 있다.

\begin{dhkcode}
obj = Example()

instance_descriptor = Example.__dict__["instance_method"]
class_descriptor = Example.__dict__["class_method"]
static_descriptor = Example.__dict__["static_method"]


a = instance_descriptor.__get__(obj, Example)
b = class_descriptor.__get__(obj, Example)
c = static_descriptor.__get__(obj, Example)
\end{dhkcode}

이제 확인해 보자.

\begin{dhkcode}
print(a.__self__ is obj)
print(b.__self__ is Example)
print(type(c))
\end{dhkcode}

실행 결과:

\begin{dhkzsh}
True
True
<class 'function'>
\end{dhkzsh}

이것이 세 메서드의 차이를 설명하는 가장 중요한 객체 모델이다.

\zsh{@classmethod}와 \zsh{@staticmethod}는 함수 본문에 특별한 능력을 부여하는 문법이 아니라, 함수가 클래스 속성으로 접근될 때의 바인딩 동작을 변경하는 디스크립터를 만드는 데코레이터다.

\subsubsection{\zsh{self}, \zsh{cls}, 정적 메서드의 선택 기준}

이제 실제로 클래스를 설계할 때 어떤 메서드를 선택해야 할지 생각해 보자.

다음 상황을 가정하자.

\begin{dhkcode}
class Person:
    species = "Human"

    def __init__(self, name, age):
        self.name = name
        self.age = age
\end{dhkcode}

세 가지 기능을 추가하고 싶다.

첫 번째는 자신의 이름을 출력하는 기능이다.

\begin{dhkcode}
def introduce(self):
    return self.name
\end{dhkcode}

이 기능은 인스턴스의 \zsh{name} 속성을 사용해야 한다.

따라서 인스턴스 메서드가 적절하다.

두 번째는 문자열로부터 인스턴스를 생성하는 기능이다.

\begin{dhkcode}
@classmethod
def from_string(cls, data):
    name, age = data.split(",")
    return cls(name, int(age))
\end{dhkcode}

이 기능은 실제 접근 대상 클래스를 사용하여 객체를 생성해야 한다.

따라서 클래스 메서드가 적절하다.

세 번째는 나이가 유효한지 확인하는 기능이다.

\begin{dhkcode}
@staticmethod
def is_valid_age(age):
    return isinstance(age, int) and age >= 0
\end{dhkcode}

이 기능은 인스턴스 상태도, 클래스 상태도 필요하지 않다.

따라서 정적 메서드로 작성할 수 있다.

최종 코드는

\begin{dhkcode}
class Person:
    species = "Human"

    def __init__(self, name, age):
        self.name = name
        self.age = age

    def introduce(self):
        return f"My name is {self.name}"

    @classmethod
    def from_string(cls, data):
        name, age = data.split(",")
        return cls(name, int(age))

    @staticmethod
    def is_valid_age(age):
        return isinstance(age, int) and age >= 0
\end{dhkcode}

이제 세 메서드는 사로 다른 목적을 갖게 된다.

\begin{dhkcode}
p = Person.from_string("Kim,30")

print(p.introduce())
print(Person.is_valid_age(p.age))
\end{dhkcode}

실행 결과:

\begin{dhkzsh}
My name is Kim
True
\end{dhkzsh}

\subsubsection{클래스 메서드와 클래스 속성 변경에서 주의할 점}

2.7에서는 클래스 속성과 인스턴스 속성의 재바인딩 차이를 학습했다.

클래스 메서드에서도 비슷한 문제가 발생할 수 있다.

다음 코드를 보자.

\begin{dhkcode}
class Counter:
    count = 0

    @classmethod
    def increase(cls):
        cls.count += 1

class ChildCounter(Counter):
    pass
\end{dhkcode}

이제 다음 코드를 실행해 보자.

\begin{dhkcode}
ChildCounter.increase()

print(Counter.count)
print(ChildCounter.count)
\end{dhkcode}

실행 결과:

\begin{dhkzsh}
0
1
\end{dhkzsh}

왜 이런 결과가 발생할까?

\zsh{ChildCounter.increase()}를 호출했기 때문에 \zsh{cls}는 \zsh{ChildCounter}다.

따라서 다음 표현식은

\begin{dhkcode}
cls.count += 1
\end{dhkcode}

개념적으로 다음 과정을 거친다.

\begin{dhksmry}
    1. ChildCounter.count를 읽는다.

    2. ChildCounter에 없으므로
        부모 Counter.count에서 0을 얻는다.

    3. 1을 더한다.

    4. ChildCounter.count에 결과를 설정한다.
\end{dhksmry}

따라서 부모 클래스의 \zsh{count}는 그대로이고, 자식 클래스에 새로운 \zsh{count} 속성이 생성된다.

\begin{dhkcode}
print(Counter.__dict__["count"])
print(ChildCounter.__dict__["count"])
\end{dhkcode}

실행 결과:

\begin{dhkzsh}
0
1
\end{dhkzsh}

이것은 2.7에서 학습했던 속성 가리기와 같은 원리다.

따라서 \zsh{@classmethod}를 사용한다고 해서 모든 클래스가 하나의 공유 상태를 자동으로 사용하는 것은 아니다.

\zsh{cls}가 어떤 클래스 객체를 참조하는지에 따라 속성 설정 위치가 달라질 수 있다.

\subsubsection{PyTorch와의 연결 - \zsh{nn.Module}에서는 어떤 메서드를 사용하는가?}

이제 PyTorch의 객체 모델과 연결해 보자.

다음 코드를 보자.

\begin{dhkcode}
import torch
import torch.nn as nn


class MyModel(nn.Module):
    def __init__(self, input_dim, output_dim):
        super().__init__()

        self.linear = nn.Linear(
            input_dim,
            output_dim
        )

    def forward(self, x):
        return self.linear(x)
\end{dhkcode}

여기서 \zsh{forward()}는 어떤 종류의 메서드일까?

인스턴스 메서드다.

\begin{dhkcode}
def forward(self, x):
\end{dhkcode}

\zsh{self}는 실제 모델 인스턴스를 참조한다.

따라서 다음 속성에 접근할 수 있다.

\begin{dhkcode}
self.linear
\end{dhkcode}

이것은 모델 인스턴스마다 등록되 ㄴ하위 모듈이기 때문에 인스턴스 메서드가 적절하다.

\begin{dhknote}
\begin{dhksmry}[centered]
클래스 메서드를 추가한다면?
\end{dhksmry}
\end{dhknote}

예를 들어 설정 딕셔너리에서 모델을 생성하는 기능을 만들 수 있다.

\begin{dhkcode}
class MyModel(nn.Module):
    def __init__(self, input_dim, output_dim):
        super().__init__()

        self.linear = nn.Linear(
            input_dim,
            output_dim
        )

    def forward(self, x):
        return self.linear(x)

    @classmethod
    def from_config(cls, config):
        return cls(
            input_dim=config["input_dim"],
            output_dim=config["output_dim"]
        )
\end{dhkcode}

이제 다음처럼 사용할 수 있다.

\begin{dhkcode}
config = {
    "input_dim": 10,
    "output_dim": 5
}

model = MyModel.from_config(config)
\end{dhkcode}

여기서 \zsh{cls}는 \zsh{MyModel} 클래스 객체다.

따라서 설정값을 이용해 모델 인스턴스를 생성할 수 있다.

이런 방식은 모델 구조를 설정 파일이나 딕셔너리로 관리할 때 유용하다.

\begin{dhknote}
\begin{dhksmry}[centered]
정적 메서드를 추가한다면?
\end{dhksmry}
\end{dhknote}

입력 차원의 유효성을 검사하는 함수를 추가할 수도 있다.

\begin{dhkcode}
class MyModel(nn.Module):
    @staticmethod
    def validate_dimension(dim):
        return isinstance(dim, int) and dim > 0
\end{dhkcode}

이 메서드는 특정 모델 인스턴스나 클래스 상태를 사용할 필요가 없다.

따라서 정적 메서드로 정의할 수 있다.

물론 검증 기능이 여러 클래스에서 공통으로 사용된다면 클래스 밖의 일반 함수로 분리하는 것도 좋은 설계다.

\subsubsection{세 메서드의 최종 비교}

지금까지 학습한 내용을 하나의 표로 정리해 보자.

\begin{dhktable}
|구분|	인스턴스 메서드|	클래스 메서드|	정적 메서드|
|:---|:---|:---|:---|
|데코레이터|	없음|	@classmethod|	@staticmethod|
|관례적 첫 인수|	self|	cls|	없음|
|자동 바인딩|	인스턴스|	클래스|	없음|
|인스턴스 상태 접근|	직접 가능|	인스턴스가 별도로 필요|	인스턴스가 별도로 필요|
|클래스 상태 접근|	type(self) 등|	cls|	클래스 이름 등을 명시|
|상속 시 바인딩|	실제 인스턴스|	접근 대상 클래스|	없음|
|대표적 용도|	객체 동작|	대체 생성자, 클래스 단위 동작|	관련 유틸리티 함수|
|클래스 __dict__|	함수 객체|	classmethod 객체|	staticmethod 객체|
\end{dhktable}

여기서 한 가지 주의할 점이 있다.

인스턴스 메서드는 클래스에서도 접근할 수 있다.

\begin{dhkcode}
Person.introduce
\end{dhkcode}

하지만 클래스에서 접근하면 인스턴스가 자동으로 바인딩되지 않으므로, 일반적으로 인스턴스를 명시적으로 전달해야 한다.

\begin{dhkcode}
Person.introduce(p)
\end{dhkcode}

반면 클래스 메서드는 클래스에서 접근해도 클래스 객체가 자동으로 바인딩된다.

\begin{dhkcode}
Person.from_string("Kim,30")
\end{dhkcode}

정적 메서드는 어느 쪽에서도 자동 바인딩이 일어나지 않는다.

\begin{dhkcode}
Person.is_valid_age(30)
p.is_valid_age(30)
\end{dhkcode}

이 차이가 세 메서드를 구분하는 가장 중요한 기준이다.

\subsubsection{핵심 정리}

이번 학습에서 반드시 이해해야 할 내용을 정리해 보자.

\begin{dhknote}
\begin{dhksmry}[centered]
세 메서드의 차이는 함수 본문보다 바인딩 방식에 있다.
\end{dhksmry}
\end{dhknote}

\begin{dhkcode}
def method(self):
    ...


@classmethod
def method(cls):
    ...


@staticmethod
def method():
    ...
\end{dhkcode}

\begin{dhknote}
\begin{dhksmry}[centered]
인스턴스 메서드는 인스턴스에 바인딩된다.
\end{dhksmry}
\end{dhknote}

\begin{dhkcode}
p.introduce.__self__ is p
\end{dhkcode}

\begin{dhknote}
\begin{dhksmry}[centered]
클래스 메서드는 클래스 객체에 바인딩된다.
\end{dhksmry}
\end{dhknote}

\begin{dhkcode}
Person.get_species.__self__ is Person
\end{dhkcode}

\begin{dhknote}
\begin{dhksmry}[centered]
정적 메서드는 자동 바인딩을 수행하지 않는다.
\end{dhksmry}
\end{dhknote}

클래스나 인스턴스로 접근해도 일반적으로 함수 객체를 반환한다.

\begin{dhknote}
\begin{dhksmry}[centered]
@classmethod와 @staticmethod는
디스크립터 객체를 생성하는 데코레이터다.
\end{dhksmry}
\end{dhknote}

클래스 네임스페이스에 저장되는 객체의 종류와 \zsh{__get__()} 동작을 변경한다.

\begin{dhknote}
\begin{dhksmry}[centered]
클래스 메서드는 상속 관계에서 유용하다.
\end{dhksmry}
\end{dhknote}

\zsh{cls}가 실제 접근 대상 클래스로 바인딩되므로, 자식 클래스의 타입을 유지하는 대체 생성자를 구현할 수 있다.

\begin{dhknote}
\begin{dhksmry}[centered]
클래스 메서드에서 클래스 속성을 변경할 때는
바인딩 대상을 고려해야 한다.
\end{dhksmry}
\end{dhknote}

\zsh{cls.count += 1}은 자식 클래스에 새로운 클래스 속성을 만들 수 있다.

\begin{dhknote}
\begin{dhksmry}[centered]
PyTorch의 forward()는 인스턴스 메서드다.
\end{dhksmry}
\end{dhknote}

모델 인스턴스의 하위 모듈과 파라미터를 사용해야 하기 때문이다.

\subsubsection{2.11에서 반드시 이해해야 할 결론}

이번 학습의 핵심은 다음과 같다.

\begin{dhktable}
|Method Type|Binding Target|
|:---|:---|
|Instance Method|Instance|
|Class Method|Class|
|Static Method|None|
\end{dhktable}

즉, Python에서 메서드의 종류를 구분하는 가장 정확한 기준은 어떤 객체가 자동으로 바인딩되는가이다.

이제 다음 코드를 보면

\begin{dhkcode}
class MyModel(nn.Module):
    def forward(self, x):
        return self.linear(x)

    @classmethod
    def from_config(cls, config):
        return cls(**config)

    @staticmethod
    def validate_dimension(dim):
        return isinstance(dim, int) and dim > 0
\end{dhkcode}

각 메서드가 왜 해당 방식으로 정의되었는지 설명할 수 있어야 한다.

\begin{dhksmry}
forward()는 모델 인스턴스의 상태를 사용한다.

from_config()는 클래스 객체를 호출하여 인스턴스를 생성한다.

validate_dimension()은 인스턴스나 클래스의 자동 바인딩이 필요하지 않다.
\end{dhksmry}

이 세 가지를 구분할 수 있다면 Python의 메서드 모델을 구조적으로 이해한 것이다.